\documentclass[a4paper, fontsize=14pt]{article}

\usepackage[T2A]{fontenc}
\usepackage[utf8]{inputenc}
\usepackage[russian]{babel}

\usepackage[unicode, pdftex]{hyperref}
\usepackage{lastpage}
\usepackage{float}
\usepackage{gensymb}
\usepackage{booktabs}
\usepackage{amsmath}
\usepackage{amssymb}
\usepackage{tikz}
\usetikzlibrary{shapes, arrows}

\usepackage{graphicx}
\usepackage{subcaption}

\usepackage{geometry}
\geometry{margin=1cm}

\title{Q-learning динамика на графе: постановка, сходимость и сравнение с динамикой Константинова}
\author{}
\date{}

% --- Работа с теоремами ---
\usepackage{amsthm} 
\usepackage[most]{tcolorbox} 

% Определяем цвета
\definecolor{theorembg}{RGB}{240, 248, 255} 
\definecolor{theoremframe}{RGB}{0, 0, 139}   

% 1. Определяем красивую теорему (заменяет стандартную theorem из amsthm)
\newtcbtheorem[number within=section]{theorem}{Теорема}%
{
    colback=theorembg,
    colframe=theoremframe,
    fonttitle=\bfseries,
    coltitle=white,
    enhanced,
    attach boxed title to top left={xshift=6mm,yshift=-3mm},
    boxed title style={colback=theoremframe},
    sharp corners,
    top=6mm
}{th}

% 2. Остальные окружения (Лемма, Определение) можно оставить стандартными 
% или тоже переделать в box. Здесь оставляем стандартными, но привязываем 
% их нумерацию к счетчику tcb-теоремы (tcb@cnt@theorem).
\theoremstyle{plain}
\newtheorem{lemma}{Лемма}[section] % Или [theorem], если удастся связать счетчики
\newtheorem{corollary}{Следствие}[section]

\theoremstyle{definition}
\newtheorem{definition}{Определение}[section]
\newtheorem{example}{Пример}[section]

\begin{document}
\maketitle

\section{Введение}
% 0.5--1 страница: мотивация...

\section{Формальная постановка модели}
\subsection{Граф и агенты}
% G=(V,E)...

\subsection{Игра и выплаты}
% A_i...

\subsection{Базовая динамика (по Константинову)}
% Строго...

\section{Q-learning динамика}
\subsection{Определение состояния и Q-функции}
% ...

\subsection{Политика выбора действий}
% ...

\subsection{Правило обновления}
% ...

\section{Сходимость}
% Синтаксис: \begin{theorem}{Заголовок}{Метка для ссылки} Текст \end{theorem}
\begin{theorem}{Предельная кооперация}{statless-coop}
    Если игра stateless, то средняя кооперация не зависит от выплаты $b$. И имеет формулу:
    \begin{enumerate}
        \item \textbf{pp-reward}: $\frac{1}{n} \sum_{i=1}^n \sigma(e^{c \cdot k_i / T})$
        \item \textbf{pf-reward}: $\frac{1}{n} \sum_{i=1}^n \sigma(e^{c / T})$
        \item \textbf{fp-reward}: $\frac{1}{n} \sum_{i=1}^n \sigma(e^{c \cdot k_i / T})$
        \item \textbf{ff-reward}: $\frac{1}{n} \sum_{i=1}^n \sigma(e^{c / T})$
    \end{enumerate}
\end{theorem}



\begin{proof}
    В начале покажем, что процесс обучения сходится к единственным значениям $Q^*$.
    Правило обновления Q-значений для агента $i$ имеет вид:
    $$ Q_{i, t+1}(a) = (1-\alpha_t) Q_{i, t}(a) + \alpha_t \left( r_{i, t} + \gamma \max_{a'} Q_{i, t}(a') \right) $$
    Это уравнение можно переписать в виде итерации стохастической аппроксимации. Поскольку $\gamma < 1$, соответствующий оператор Беллмана является сжимающим. При выполнении стандартных условий Роббинса-Монро ($\sum \alpha_t = \infty, \sum \alpha_t^2 < \infty$), значения $Q_i(a)$ с вероятностью 1 сходятся к неподвижной точке $Q_i^*(a)$.

    Рассмотрим уравнение Беллмана, определяющее значения Q-функции в стационарном состоянии:
    $$ Q_i^*(a) = \mathbb{E}[R_i(a)] + \gamma \max_{a'} Q_i^*(a') $$
    При больцмановской политике вероятность выбора действия определяется как:
    $$ P_i(a) = \frac{e^{Q_i(a)/T}}{\sum_{a'} e^{Q_i(a')/T}} $$
    Вычислим разность Q-значений для действий кооперации (C) и отказа (D). Заметим, что слагаемое с максимумом по будущим действиям сокращается:
    $$ \Delta Q_i = Q_i(C) - Q_i(D) = \left[R_i(C) + \gamma \max_{a'} Q_i^*(a')\right] - \left[R_i(D) + \gamma \max_{a'} Q_i^*(a')\right] = R_i(C) - R_i(D)$$
    В среднем:
    $$ \Delta Q_i^* =  \mathbb{E}[R_i(C)] - \mathbb{E}[R_i(D)]$$
    
    Выразим вероятность кооперации через сигмоидную функцию:
    $$ P_i(C) = \sigma\left(-\frac{\Delta Q_i}{T}\right)$$
    $$ P_i^*(C) = \sigma\left(-\frac{\Delta Q_i^*}{T}\right)$$
    
    Тогда средний уровень кооперации в сети $\hat{\rho}$ равен:
    $$\hat{\rho} = \frac{1}{N} \cdot \left( \sum_{i=0}^N\mathbb{E}x_i\right) = \frac{1}{N} \cdot \left( \sum_{i=0}^N P_i^*(C)\right) = \frac{1}{N} \cdot \left( \sum_{i=0}^N \sigma\left(-\frac{\Delta Q_i^*}{T}\right)\right)$$
    
    Теперь рассмотрим каждый тип функции выплат отдельно. Заметим, что выигрыш от действий соседей $b$ сокращается при вычислении разности Q-значений.
    
    1. В случае награды pp (pay-per-partner / pay-per-partner):
    $$ \mathbb{E}[R_i(a)] = b \cdot \sum_{j \in N_i} \mathbb{E}[x_j] - c \cdot k_i \cdot a $$
    Разность Q-значений:
    $$ \Delta Q_i^* = (b \cdot \sum_{j \in N_i} \mathbb{E}[x_j] - c \cdot k_i \cdot 1)  - (b \cdot \sum_{j \in N_i} \mathbb{E}[x_j] - c \cdot k_i \cdot 0) = - c \cdot k_i $$
    Итоговая формула для средней кооперации:
    $$\hat{\rho} = \frac{1}{N} \cdot \left( \sum_{i=0}^N \sigma\left(\frac{c \cdot k_i}{T} \right)\right)$$ 
    
    
    2. В случае награды pf (pay-per-partner / fixed-cost):
    $$ \mathbb{E}[R_i(a)] = b \cdot \sum_{j \in N_i} \mathbb{E}[x_j] - c \cdot a $$
    Разность Q-значений:
    $$ \Delta Q_i^* = (b \cdot \sum_{j \in N_i} \mathbb{E}[x_j] - c \cdot 1)  - (b \cdot \sum_{j \in N_i} \mathbb{E}[x_j] - c \cdot 0) = - c $$
    Итоговая формула:
    $$\hat{\rho} = \frac{1}{N} \cdot \left( \sum_{i=0}^N \sigma\left(\frac{c}{T} \right)\right)$$
    
    
    3. В случае награды fp (fixed-benefit / pay-per-partner):
    $$ \mathbb{E}[R_i(a)] = b \cdot \sum_{j \in N_i} \frac{\mathbb{E}[x_j]}{k_j} - c \cdot k_i \cdot a $$
    Разность Q-значений:
    $$ \Delta Q_i^* = (b \cdot \sum_{j \in N_i} \frac{\mathbb{E}[x_j]}{k_j} - c \cdot k_i \cdot 1)  - (b \cdot \sum_{j \in N_i} \frac{\mathbb{E}[x_j]}{k_j} - c \cdot k_i \cdot 0) = - c \cdot k_i $$
    Итоговая формула:
    $$\hat{\rho} = \frac{1}{N} \cdot \left( \sum_{i=0}^N \sigma\left(\frac{c \cdot k_i}{T} \right)\right)$$
    
    
    4. В случае награды ff (fixed-benefit / fixed-cost):
    $$ \mathbb{E}[R_i(a)] = b \cdot \sum_{j \in N_i} \frac{\mathbb{E}[x_j]}{k_j} - c \cdot a $$
    Разность Q-значений:
    $$ \Delta Q_i^* = (b \cdot \sum_{j \in N_i} \frac{\mathbb{E}[x_j]}{k_j} - c \cdot 1)  - (b \cdot \sum_{j \in N_i} \frac{\mathbb{E}[x_j]}{k_j} - c \cdot 0) = - c $$
    Итоговая формула:
    $$\hat{\rho} = \frac{1}{N} \cdot \left( \sum_{i=0}^N \sigma\left(\frac{c}{T} \right)\right)$$
\end{proof}


\subsection{Обсуждение и сравнение с динамикой Константинова}
% ...

\section{Численные эксперименты}
\subsection{Настройки эксперимента}
% ...

\subsection{Графики сходимости}
В данном разделе представлены результаты симуляций для различных типов вознаграждения ($pp$, $pf$, $fp$, $ff$) и различных значений параметра $b$. Для каждого типа вознаграждения приведены графики сходимости при разных значениях коэффициента дисконтирования $\gamma$.

\begin{figure}[H]
    \centering
    \begin{subfigure}{0.48\textwidth}
        \includegraphics[width=\textwidth]{pp_b1.5_gamma_comparison.png}
        \caption{$b=1.5$}
    \end{subfigure}
    \hfill
    \begin{subfigure}{0.48\textwidth}
        \includegraphics[width=\textwidth]{pp_b2.5_gamma_comparison.png}
        \caption{$b=2.5$}
    \end{subfigure}
    \begin{subfigure}{0.48\textwidth}
        \includegraphics[width=\textwidth]{pp_b3.5_gamma_comparison.png}
        \caption{$b=3.5$}
    \end{subfigure}
    \caption{Сходимость для PP-reward}
\end{figure}

\begin{figure}[H]
    \centering
    \begin{subfigure}{0.48\textwidth}
        \includegraphics[width=\textwidth]{pf_b1.5_gamma_comparison.png}
        \caption{$b=1.5$}
    \end{subfigure}
    \hfill
    \begin{subfigure}{0.48\textwidth}
        \includegraphics[width=\textwidth]{pf_b2.5_gamma_comparison.png}
        \caption{$b=2.5$}
    \end{subfigure}
    \vfill
    \begin{subfigure}{0.48\textwidth}
        \includegraphics[width=\textwidth]{pf_b3.5_gamma_comparison.png}
        \caption{$b=3.5$}
    \end{subfigure}
    \caption{Сходимость для PF-reward}
\end{figure}

\begin{figure}[H]
    \centering
    \begin{subfigure}{0.48\textwidth}
        \includegraphics[width=\textwidth]{fp_b1.5_gamma_comparison.png}
        \caption{$b=1.5$}
    \end{subfigure}
    \hfill
    \begin{subfigure}{0.48\textwidth}
        \includegraphics[width=\textwidth]{fp_b2.5_gamma_comparison.png}
        \caption{$b=2.5$}
    \end{subfigure}
    \begin{subfigure}{0.48\textwidth}
        \includegraphics[width=\textwidth]{fp_b3.5_gamma_comparison.png}
        \caption{$b=3.5$}
    \end{subfigure}
    \caption{Сходимость для FP-reward}
\end{figure}

\begin{figure}[H]
    \centering
    \begin{subfigure}{0.48\textwidth}
        \includegraphics[width=\textwidth]{ff_b1.5_gamma_comparison.png}
        \caption{$b=1.5$}
    \end{subfigure}
    \hfill
    \begin{subfigure}{0.48\textwidth}
        \includegraphics[width=\textwidth]{ff_b2.5_gamma_comparison.png}
        \caption{$b=2.5$}
    \end{subfigure}
    \begin{subfigure}{0.48\textwidth}
        \includegraphics[width=\textwidth]{ff_b3.5_gamma_comparison.png}
        \caption{$b=3.5$}
    \end{subfigure}
    \caption{Сходимость для FF-reward}
\end{figure}

\section{Заключение}
% ...

% \bibliographystyle{plain}
% \bibliography{refs}

\end{document}
